\section{投资问题：高利润能否变成可分配现金}

2025年是低谷，2026Q1和H1预告确认利润反转；但估值能否扩张，取决于利润是否覆盖营运资金、利息和资本开支。核心判断是，毛利率和归母利润改善显著，Q1经营现金流却为-25.73亿元，存货较年初增加53.94亿元，债务代理和在建项目仍高。现阶段应把利润预测与估值倍数分开：利润反转支持77.5亿元基准，现金缺口限制基准PE为7.5倍。

\begin{exhibitbox}[财务质量总览]
\noindent\begin{tabularx}{\textwidth}{L{3.25cm}R{2.1cm}R{2.1cm}X}
\toprule
\textbf{指标} & \textbf{2025A} & \textbf{2026Q1} & \textbf{投资解读} \\
\midrule
营业收入 & 1,135.27亿元 & 299.48亿元 & Q1收入规模恢复 \\
归母净利润 & 2.58亿元 & 19.95亿元 & 利润反转显著 \\
扣非归母 & 1.41亿元 & 19.40亿元 & Q1改善以核心经营为主 \\
毛利率 & 4.51\% & 12.88\% & 产品价差与负荷改善 \\
少数股东损益 & 0.37亿元 & 8.23亿元 & 高景气归属折损显著 \\
经营现金流 & 46.25亿元 & -25.73亿元 & 利润尚未同步变现 \\
存货 & 133.13亿元 & 187.07亿元 & 一季度增加53.94亿元 \\
在建工程 & 149.57亿元 & 136.95亿元 & 资本占用仍高 \\
购建长期资产现金支出 & 55.02亿元 & 14.43亿元 & 自由现金流受资本开支约束 \\
资产负债率 & 72.28\% & 71.46\% & 利润增厚权益但杠杆仍高 \\
\bottomrule
\end{tabularx}
\end{exhibitbox}
\sourcenote{HY-OFF-2025AR（已审计）；HY-OFF-2026Q1（未经审计季报）；data/verified\_financials.md，财务截止2026-03-31。毛利率为收入减营业成本派生。}

利润率从4.51\%升至12.88\%，与文莱油品／芳烃高景气方向一致；扣非接近归母，说明Q1并非主要依赖资产处置或公允价值。但同一季度存货和债务上升，现金流大额流出，表明原料采购、备货、产品库存和结算时点可能吸收了利润。H1现金流能否反转，是中报最重要的质量指标。

\section{2025低谷：利润薄，但现金并未完全失守}

2025年收入1,135.27亿元、归母2.58亿元，净利率极低；经营现金流46.25亿元，高于归母，却仍低于购建长期资产现金支出55.02亿元，简单自由现金流代理约为-8.76亿元。公司财务费用28.11亿元，其中利息费用23.36亿元，远高于当年归母利润，说明低景气时债务成本对普通股东利润具有强放大效应。

2025年信用与资产减值合计减少利润总额约0.41亿元，其中存货跌价损失约0.49亿元；考虑所得税和少数股东后减少归母约0.31亿元。减值不是2025低利润的主要原因，核心仍是分部毛利低与财务负担。若2026商品价格快速回落，库存减值风险可能重新放大，但不应预先把全部库存增量视为损失。

\begin{exhibitbox}[2025利润到现金与资本回报]
\noindent\begin{tabularx}{\textwidth}{L{4.1cm}R{2.2cm}X}
\toprule
\textbf{项目} & \textbf{2025A} & \textbf{含义} \\
\midrule
归母净利润 & 2.58亿元 & 周期低谷的股东利润 \\
经营现金流 & 46.25亿元 & 高于净利，受折旧和营运资金影响 \\
购建长期资产现金支出 & 55.02亿元 & 超过OCF \\
简单自由现金流代理 & 约-8.76亿元 & 不含资产处置和其他投资现金流 \\
财务费用／利息费用 & 28.11／23.36亿元 & 低景气时吞噬归母 \\
减值对归母影响 & -0.31亿元 & 不是低利润主因 \\
\bottomrule
\end{tabularx}
\end{exhibitbox}
\sourcenote{HY-OFF-2025AR；2025资产减值公告HY-OFF-2025AR配套文件；现金流代理为OCF减购建长期资产支出，未等同公司披露FCF。}

投资含义是，2026利润提升必须先补偿资本成本和现金缺口，才能形成可分配价值。2025已经说明，低景气时仅靠产业规模无法覆盖利息和资本开支；因此不宜用高峰ROE给出成长股PB。

\section{Q1现金缺口：库存、采购与结算共同作用}

2026Q1存货由133.13亿元升至187.07亿元，增幅约40.5\%；经营现金流-25.73亿元，购建长期资产支出14.43亿元，简单自由现金流代理约-40.17亿元。与此同时，归母利润19.95亿元。利润与现金之间约45.68亿元的方向差，说明营运资金是首要解释变量。

库存上升可能来自油价与原料价格、文莱和国内备货、产成品库存及在途物资；Q1报告只披露总额与原因，未提供分产品库存和跌价。投资者不应把库存全部视为风险，也不能把它自动视为未来利润。半年报需要核验库存是否随销售转化、经营现金流是否转正以及短债是否下降。

\begin{exhibitbox}[Q1利润—现金缺口桥]
\noindent\begin{tabularx}{\textwidth}{L{4.0cm}R{2.2cm}X}
\toprule
\textbf{项目} & \textbf{2026Q1} & \textbf{观察} \\
\midrule
归母净利润 & +19.95亿元 & 会计利润显著改善 \\
经营现金流 & -25.73亿元 & 与归母方向相反 \\
利润与OCF方向差 & -45.68亿元 & 营运资金、税息和少数股东需拆解 \\
存货增加 & +53.94亿元 & 是现金占用重要候选 \\
购建长期资产现金支出 & -14.43亿元 & 项目与技改继续占用 \\
简单自由现金流代理 & -40.17亿元 & OCF减购建支出，非公司正式FCF \\
\bottomrule
\end{tabularx}
\end{exhibitbox}
\sourcenote{HY-OFF-2026Q1；data/verified\_financials.md，财务截止2026-03-31；方向差与自由现金流代理为AStock派生。}

\begin{keyinsight}[现金流的“所以呢”]
利润预测和估值倍数必须分开更新。\textbf{行动上，若H1归母落在预告区间但OCF仍大额为负，本报告不必立即下调77.5亿元利润，却应把7.5倍PE向6倍收敛；若现金转正并覆盖利息和资本开支，才允许提高倍数。}
\end{keyinsight}

\section{债务与流动性：利润上升不等于杠杆问题结束}

2025年末短期借款396.96亿元、长期借款130.83亿元、一年内到期非流动负债95.72亿元、应付债券27.44亿元，货币资金107.49亿元。以这些主要有息负债减货币资金的粗略净债务代理约543.46亿元；其中约40.56亿元货币资金使用受限，因此可自由动用现金低于账面货币资金。

2026Q1上述四项负债分别为427.69、167.10、77.74和27.71亿元，货币资金126.26亿元，粗略净债务代理约573.98亿元，较年末增加约30.52亿元。该代理未包含全部租赁、票据、衍生与受限现金调整，不是审计净债务，但足以显示一季度利润改善没有同步减少融资敞口。

\begin{exhibitbox}[主要债务与现金桥]
\small
\noindent\begin{tabularx}{\textwidth}{L{3.4cm}R{2.05cm}R{2.05cm}X}
\toprule
\textbf{项目} & \textbf{2025年末} & \textbf{2026Q1} & \textbf{变化与含义} \\
\midrule
短期借款 & 396.96亿元 & 427.69亿元 & 增加30.73亿元 \\
一年内到期非流动负债 & 95.72亿元 & 77.74亿元 & 到期结构下降17.98亿元 \\
长期借款 & 130.83亿元 & 167.10亿元 & 增加36.27亿元 \\
应付债券 & 27.44亿元 & 27.71亿元 & 含实际利率法影响 \\
货币资金 & 107.49亿元 & 126.26亿元 & 账面现金增加18.77亿元 \\
粗略净债务代理 & 543.46亿元 & 573.98亿元 & 增加约30.52亿元 \\
资产负债率 & 72.28\% & 71.46\% & 权益因利润增加，绝对债务仍高 \\
\bottomrule
\end{tabularx}
\end{exhibitbox}
\sourcenote{HY-OFF-2025AR；HY-OFF-2026Q1。粗略净债务代理为短借+一年内到期非流动负债+长借+应付债券-货币资金，未调整全部受限现金和其他融资。}

流动性风险不等于立即偿债危机。公司有银行授信、产业资产和高景气利润支持再融资；但短债和资本开支同时存在会限制分红、回购和估值扩张。文莱二期若继续提高长期借款，必须用有约束力融资和项目现金流匹配，而不是依赖峰值利润滚动补充资本金。

\section{少数股东与归母转换：Q1已经给出警告}

2026Q1合并净利润约28.18亿元，其中归母19.95亿元、少数股东损益8.23亿元。少数股东约占29.2\%，与文莱70\%权益方向接近。若文莱是主要增量利润池，高景气越强，少数股东绝对分流越大；这不是一次性损失，而是结构性归属。

半年报应同时披露总净利润、归母、少数股东损益和主要非全资子公司财务。若少数股东比例显著高于30\%且没有其他合营／子公司解释，House从项目利润向归母的转换率应下调。反之，若国内全资／高持股资产利润改善，归母占比可能上升。

\section{财务质量结论与更新规则}

2026盈利反转已经成立，现金质量仍为未完成验证。中报升级条件是：归母和扣非落在预告区间、经营现金流较Q1显著改善并最好转正、存货下降或增速显著低于收入、短债不再上升、少数股东归属可解释。任何一项单独改善都不足以解除资本折价。

\paragraph{财务失效条件} 若H1经营现金流仍大额为负、存货和短期借款继续同步上升、利息费用未因转股下降，或少数股东分流显著高于预期，则高利润没有形成相应每股现金价值。\textbf{研究上，下调估值倍数并把16.5元以上设为重新评估区。}

相反，若H1 OCF转正并覆盖利息和资本开支，且净债务代理下降，现金质量将成为上调倍数而非上调利润的依据。保持这一分工，可避免同一利好同时抬升利润和PE，造成双重计算。
